\documentclass[12pt]{article}

\usepackage[a4paper, margin=1in]{geometry}
\usepackage{graphicx}
\usepackage{hyperref}
\usepackage{booktabs}
\usepackage{float}
\usepackage{xcolor}
\usepackage{setspace}
\usepackage{titlesec}
\usepackage{enumitem}
\usepackage{times}   % Times font for pdfLaTeX

\onehalfspacing

\titleformat{\section}{\large\bfseries}{\thesection}{1em}{}
\titleformat{\subsection}{\normalsize\bfseries}{\thesubsection}{1em}{}

\title{\textbf{Mini Project 1: Exploratory Data Analysis (EDA)}}
\author{Achinthya M \\ Data Science \& AI Internship}
\date{}

\begin{document}
	
	\maketitle
	\thispagestyle{empty}
	
	\tableofcontents
	\newpage
	
	% --------------------------------------------------
	\section{Introduction}
	
	This mini project focuses on performing Exploratory Data Analysis (EDA) on a real-world dataset. The objective was to understand the structure of the dataset, clean missing values, visualize patterns, and extract meaningful insights.
	
	The dataset analyzed in this project is \textbf{customer\_analytics.csv}, where each row represents an individual customer with demographic and behavioral attributes.
	
	% --------------------------------------------------
	\section{Project Structure}
	
	The project follows the directory layout below:
	
	\begin{flushleft}
		\texttt{
			MiniProject1/\\
			\quad data/\\
			\quad\quad customer\_analytics.csv\\
			\quad\quad customer\_analytics\_cleaned.csv\\
			\quad notebooks/\\
			\quad\quad MiniProject1\_EDA.ipynb\\
			\quad reports/\\
			\quad\quad MiniProject1\_Documentation.tex\\
			\quad README.md\\
			\quad requirements.txt
		}
	\end{flushleft}
	
	% --------------------------------------------------
	\section{Files Description}
	
	\begin{itemize}[leftmargin=1.5cm]
		\item \textbf{customer\_analytics.csv} – Original dataset.
		\item \textbf{customer\_analytics\_cleaned.csv} – Cleaned dataset after preprocessing.
		\item \textbf{MiniProject1\_EDA.ipynb} – Notebook containing the complete EDA workflow.
		\item \textbf{README.md} – Project overview and instructions.
		\item \textbf{requirements.txt} – Required Python libraries.
	\end{itemize}
	
	% --------------------------------------------------
	\section{Phase 1: Data Inspection}
	
	The dataset was loaded using Pandas. Initial inspection was performed using:
	
	\begin{itemize}[leftmargin=1.5cm]
		\item \texttt{df.head()}
		\item \texttt{df.info()}
		\item \texttt{df.describe()}
	\end{itemize}
	
	This helped identify the structure, data types, and presence of missing values.
	
	% --------------------------------------------------
	\section{Phase 2: Data Cleaning and Preprocessing}
	
	Missing numerical values were filled using the median of each column to reduce the impact of outliers. Duplicate records were identified and removed to maintain data integrity.
	
	The cleaned dataset was saved as:
	
	\texttt{customer\_analytics\_cleaned.csv}
	
	% --------------------------------------------------
	\section{Phase 3: Univariate and Bivariate Analysis}
	
	Histograms and boxplots were created to understand variable distributions and detect outliers. Scatter plots were used to study relationships between variables such as age and income.
	
	% --------------------------------------------------
	\section{Phase 4: Correlation Analysis}
	
	A correlation matrix was generated using \texttt{df.corr()} and visualized using a heatmap to identify strong relationships between numerical features.
	
	% --------------------------------------------------
	\section{Key Insights}
	
	\begin{enumerate}[leftmargin=1.5cm]
		\item Moderate correlations were observed between certain numerical variables.
		\item Missing values required careful handling to avoid bias.
		\item Outliers were identified in selected numerical columns.
	\end{enumerate}
	
	% --------------------------------------------------
	\section{Conclusion}
	
	This project successfully implemented a structured EDA workflow. The dataset was inspected, cleaned, visualized, and analyzed systematically to extract meaningful insights.
	
	The final outputs include:
	\begin{itemize}[leftmargin=1.5cm]
		\item Cleaned dataset
		\item Jupyter Notebook with complete analysis
		\item Technical documentation report
	\end{itemize}
	
\end{document}